\section{A three-clause proof using prefix exposure}\label{sec:appendix-c}
\label{sec:prefix-exposure}

We give a second variance estimate for three-clause formulas, revealing
clause types only through the current prefix.  The argument truncates the
number of clauses testing each variable of an omitted clause at
\(\lceil\log n\rceil\), and pays separately for clauses meeting two or
more of those variables.  It gives an \(O(n\log^4 n)\) bound.  The
cube of the incidence cutoff contributes three logarithmic factors;
the critical occupation integral contributes the fourth.  The main
argument in Sections~\ref{sec:deletion}--\ref{sec:occupation} instead
averages the bounded third moment of the actual incidence count, giving
\(O(n\log n)\).  Comparing the two proofs shows where that improvement
enters.  The replacement step below also gives a distinct finite-variable
estimate: replacing one two-clause on \(u\) variables by \(g\)
three-clauses loses at most \(64(g^{-2}+u^{-2})\) in killing probability.

Throughout this section, \(C_1,C_2,\ldots\) are independent uniform
proper three-clauses on \(n\ge3\) variables: a clause uses three distinct
variables and independent fair signs, while whole clauses may repeat.
Write \(H_n\) for the index of the first unsatisfiable prefix and put
\[
 L=10n,\qquad T_n=\min\{H_n,L\},\qquad m_n=\EE T_n.
\]

\begin{theorem}[Variance from prefix exposure]\label{thm:prefix-variance}
There is an absolute constant \(C\) such that, for all integers \(n\ge3\),
\[
 \Var(T_n)\le Cn(\log n)^4.
\]
Consequently, for every \(t>0\),
\[
 \PP\bigl(|T_n-m_n|\ge t\sqrt n(\log n)^2\bigr)\le Ct^{-2}.
\]
\end{theorem}

The proof has two uses of the same lifetime estimate.  It bounds the
remaining duration of a deletion event and the total time spent at a
given level of the probability of killing the surviving assignments.
Three independent root-flip tests connect those two bounds.

\subsection{Replacing pairs by triples on a finite variable set}

For an integer \(u\ge3\) and \(S\subseteq\{0,1\}^u\), a collection of clauses
\emph{kills} \(S\) if no member of \(S\) satisfies all of them.  A
uniform proper \(r\)-clause here is chosen from the
\(2^r\binom ur\) signed clauses on \(r\) distinct variables.  For an
integer \(d\ge1\), let \(q_d(S)\) be the probability that \(d\)
independent uniform proper two-clauses kill \(S\).  Let \(\tau(S)\)
be the expected length of the first prefix of independent uniform proper
three-clauses whose conjunction kills \(S\), and set
\(\tau(\varnothing)=0\).  For nonempty \(S\), this
expectation is finite because the probability of survival after \(m\)
clauses is at most \(|S|(7/8)^m\).

Replacing a shorter constraint by native clauses has a qualitative
predecessor in Friedgut's half-cube argument
\cite[Lemma~5.7]{Friedgut}.  The residual-set conditioning used below
is the averaged replacement method of Carenini
\cite[Theorem~4.3, Lemma~6.1, and Corollary~6.2]{Carenini}; the following
bound supplies the finite-variable pair-to-triple estimate needed here.

\begin{lemma}[One-pair replacement]\label{lem:prefix-one-pair}
For every integer \(u\ge3\), every \(S\subseteq\{0,1\}^u\), and every integer
\(g\ge1\), if the \(g\) triples are independent uniform proper
three-clauses on these \(u\) variables, then
\begin{equation}\label{eq:prefix-one-pair}
 \PP(\text{the \(g\) triples kill \(S\)})
 \ge q_1(S)-64(g^{-2}+u^{-2}).
\end{equation}
\end{lemma}

\begin{proof}
If \(S=\varnothing\), both killing probabilities are one.  Suppose that
\(S\ne\varnothing\), and let \(b\) be the number of variables taking a
common value on all members of \(S\).  A single clause kills \(S\)
exactly when every one of its variables is among these \(b\) frozen
variables and every sign opposes the common value.  Thus the killing
probability of a proper \(r\)-clause is
\[
 h_r=\frac{(b)_r}{2^r(u)_r},\qquad r=2,3,
\]
where \((x)_r=x(x-1)\cdots(x-r+1)\), and \((b)_r=0\) for integers
\(0\le b<r\).  For \(b\ge3\),
\[
 \frac{h_3}{h_2^{3/2}}
 =\frac{b-2}{\sqrt{b(b-1)}}\frac{\sqrt{u(u-1)}}{u-2}
 \ge\frac14.
\]
Indeed, the first factor is at least \((b-2)/b\ge1/3\), and the
second is at least one.

Put \(h=h_2\le1/4\).  The block kills \(S\) whenever one of its
clauses kills \(S\) by itself, so its killing probability is at least
\(1-(1-h_3)^g\ge1-e^{-gh_3}\).  If \(b\ge3\) and
\((g/4)\sqrt h\ge2\), then \(gh_3\ge2h\), and this probability is at
least \(1-e^{-2h}\ge h\).  The last inequality follows because
\(1-e^{-2h}-h\) vanishes at zero and has nonnegative derivative on
\([0,1/4]\).  In the remaining case with \(b\ge3\), \(h\le64g^{-2}\).
If \(b<3\), then
\[
 h\le\frac1{2u(u-1)}\le u^{-2}.
\]
The nonnegativity of the block killing probability now proves
\eqref{eq:prefix-one-pair} in every case.
\end{proof}

\begin{remark}
For \(k=3\), the cutoff-free estimate in Lemma~\ref{lem:one-replacement}
has error \(256g^{-2}\).  Neither that error nor
\(64(g^{-2}+u^{-2})\) is uniformly smaller: the latter is smaller for
\(g<\sqrt3\,u\), the former is smaller for \(g>\sqrt3\,u\), and they
are equal when \(g=\sqrt3\,u\).  The two estimates therefore retain
different quantitative information.
\end{remark}

\begin{proposition}[Lifetime above the finite-variable cutoff]
\label{prop:prefix-lifetime}
Let \(K=64\).  There is an absolute constant \(K_1\) such that for
every integer \(u\ge3\), every nonempty \(S\subseteq\{0,1\}^u\), every integer
\(d\ge1\), and every \(a\) with \(4Kd/u^2\le a\le1\),
\begin{equation}\label{eq:prefix-lifetime}
 q_d(S)\ge a\quad\Longrightarrow\quad
 \tau(S)\le K_1d^{3/2}a^{-3/2}.
\end{equation}
\end{proposition}

\begin{proof}
Replace \(d\) independent two-clauses one at a time by independent
blocks of \(g\) proper three-clauses.  At each step condition on the
clauses in every other position, including the blocks already inserted.
Their surviving assignments form a set \(S'\subseteq S\).  The pair
being removed and the new block are independent of this conditioning.
Lemma~\ref{lem:prefix-one-pair}, including its empty-set case, bounds
the conditional loss by \(K(g^{-2}+u^{-2})\).  Averaging and telescoping
over the replacements gives
\begin{equation}\label{eq:prefix-many-pairs}
 \PP(\text{\(dg\) independent uniform proper three-clauses kill \(S\)})
 \ge q_d(S)-Kd(g^{-2}+u^{-2}).
\end{equation}
Choose \(g=\lceil\sqrt{4Kd/a}\rceil\).  Each of \(Kd/g^2\) and
\(Kd/u^2\) is at most \(a/4\), so a block of \(dg\) triples kills
the original set \(S\) with probability at least \(a/2\).  In an
infinite independent sequence, disjoint blocks have independent events
of killing that original set.  Their first successful block has expected
index at most \(2/a\), and the cumulative sequence kills \(S\) by
the end of that block.  Therefore
\[
 \tau(S)\le\frac{2dg}{a}\le 8\sqrt K\,d^{3/2}a^{-3/2},
\]
where \(g\le2\sqrt{4Kd/a}\) follows from \(d\ge1\) and \(a\le1\).
This proves the assertion, for example with \(K_1=64\).
\end{proof}

\subsection{A deleted clause and an unexposed future}

We now apply the lifetime estimate to the delay caused by omitting one
clause.  We first assume \(n\) is sufficiently large, in particular
\(n\ge6\); the final constant will cover the remaining sizes.  Keep
the original clause indices, and for \(1\le i\le L\) and
integer \(m\ge0\) put
\[
 B_m^{(i)}=\bigwedge_{\substack{1\le j\le m\\j\ne i}}C_j,
 \qquad
 T^{(i)}=\min\bigl\{\inf\{m\in\mathbb Z_{\ge0}:
                   B_m^{(i)}\notin\sat\},L\bigr\},
 \qquad D_i=T^{(i)}-T_n,
\]
with \(\inf\varnothing=\infty\).  Then \(D_i\ge0\), and
\(T^{(i)}\) does not use coordinate \(C_i\).  The coordinate-omission
inequality, Lemma~\ref{lem:coordinate-omission}, gives
\begin{equation}\label{eq:prefix-omission}
 \Var(T_n)\le\sum_{i=1}^{L}\EE D_i^2.
\end{equation}
Its martingale proof applies to precisely these index-preserving
deletions and to the cap \(L=10n\).

Fix \(i\) and condition on any particular clause \(C_i\).  Until the
end of the variance proof, probabilities refer to this conditional law;
the other clauses remain independent and uniform.  We call indices other
than \(i\) the baseline indices for this deletion.  Let \(R\) be the
three variables of \(C_i\), called its \emph{roots}, and let
\(t\in\{0,1\}^{R}\) be their unique values that falsify \(C_i\).
Put \(u=n-3\), abbreviate \(B_m=B_m^{(i)}\), and let
\(\mathcal F_m=\sigma(C_j:1\le j\le m,\ j\ne i)\) under the
conditional law.  In particular, \(u\ge3\).

For \(i\le m<L\), define
\begin{equation}\label{eq:prefix-delay-events}
 A_m=\{T_n\le m<T^{(i)}\}
     =\{B_m\in\sat,\ B_m\wedge C_i\notin\sat\}.
\end{equation}
On \(A_m\), every solution of \(B_m\) has root values \(t\).
The two processes agree before index \(i\), so
\begin{equation}\label{eq:prefix-delay-sum}
 D_i=\sum_{m=i}^{L-1}\ind{A_m}.
\end{equation}
In particular \(D_L=0\).  For every \(m\ge i\), let
\(S_m\subseteq\{0,1\}^u\) be the nonroot restrictions of solutions
of \(B_m\) with root values \(t\).  The sets are decreasing,
\(\mathcal F_m\)-measurable, and nonempty on \(A_m\).

A clause is \emph{ordinary} if it uses only nonroots.  A fresh clause
is ordinary with probability
\[
 p_0=\frac{\binom{u}{3}}{\binom n3},
\]
which is bounded below by a positive absolute constant for sufficiently
large \(n\).  Let \(U_m\) be the number of trials after index \(m\),
including the last trial, until the ordinary clauses among them kill
the original set \(S_m\); put \(U_m=0\) when \(S_m=\varnothing\).
These trials may continue past \(L\).  Conditional on \(\mathcal F_m\),
the future ordinary clauses are independent uniform proper triples on
the nonroots.  Their arrival gaps, counting the arrival trial, are
independent geometric variables of mean \(1/p_0\), independent of the
ordinary clauses themselves.  Conditioning first on the latter clauses
and summing the means of the gaps through their killing time proves
\begin{equation}\label{eq:prefix-future-mean}
 \EE[U_m\mid\mathcal F_m]=\frac{\tau(S_m)}{p_0}.
\end{equation}
The assertion uses the fact that only the prefix has been revealed;
for \(m\ge i\), all trials after \(m\) remain fresh even after fixing
\(C_i\).

\begin{lemma}[Remaining deletion time]\label{lem:prefix-future-delay}
Set \(\mathcal W_m=\min\{L,\tau(S_m)/p_0\}\).  For \(i\le m<L\),
on \(A_m\),
\[
 \EE\left[\sum_{\ell=m+1}^{L-1}\ind{A_\ell}
                  \,\middle|\,\mathcal F_m\right]\le\mathcal W_m.
\]
Consequently,
\begin{equation}\label{eq:prefix-weighted-delay}
 \EE D_i^2\le3\sum_{m=i}^{L-1}\EE[\ind{A_m}\mathcal W_m].
\end{equation}
\end{lemma}

\begin{proof}
On \(A_m\), any later solution of \(B_\ell\) still has root values
\(t\), and its nonroot part lies in \(S_m\).  Once the subsequent
ordinary clauses kill \(S_m\), no such solution remains.  Thus
\(A_\ell\) is impossible for \(\ell\ge m+U_m\).  The number of later
indices is at most both \(U_m\) and \(L\); its conditional expectation
is therefore at most \(\min\{\EE[U_m\mid\mathcal F_m],L\}\).
Equation~\eqref{eq:prefix-future-mean} proves the first assertion.

On \(A_m\), the nonempty set \(S_m\) requires at least one triple to
be killed, so \(\mathcal W_m\ge1\).  Expanding
\eqref{eq:prefix-delay-sum}, conditioning each cross term on its earlier
history, and using the first assertion gives
\begin{align*}
 \EE D_i^2
 &=\sum_m\EE\ind{A_m}
   +2\sum_m\EE\left[\ind{A_m}
      \EE\left[\sum_{\ell=m+1}^{L-1}\ind{A_\ell}
                  \,\middle|\,\mathcal F_m\right]\right]\\
 &\le\sum_m\EE[\ind{A_m}(1+2\mathcal W_m)]
 \le3\sum_m\EE[\ind{A_m}\mathcal W_m],
\end{align*}
where every outer sum runs from \(m=i\) to \(L-1\).
\end{proof}

\subsection{The three root-flip tests}

The weight in \eqref{eq:prefix-weighted-delay} is determined by the
surviving set.  To estimate the accompanying deletion probability, we
hide the parts of clauses that are already satisfied at the fixed root
values.  They will test whether a root can be flipped.

The \emph{type} of a baseline clause through index \(m\) records the
roots it contains and the signs of its root literals.  For \(j\in R\),
let \(\mathcal I_{j,m}\) be the indices other than \(i\), at most
\(m\), whose clause contains exactly root \(j\) and whose root literal
is true at \(t\); put \(s_{j,m}=|\mathcal I_{j,m}|\).  Call a clause
containing at least two roots a \emph{collision}, and let \(v_m\) count
such clauses at baseline indices through \(m\).  Finally, let
\(J_m\subseteq R\) contain those roots \(j\) for which no collision through \(m\)
has exactly one root literal true under \(t\), at root \(j\).  A collision
can exclude at most one root, hence
\begin{equation}\label{eq:prefix-available-roots}
 |J_m|\ge3-v_m.
\end{equation}

Let \(\mathcal G_m\) reveal all these prefix types and the entire signed
nonroot part of every baseline clause through \(m\), except at indices
in \(\bigcup_{j\in R}\mathcal I_{j,m}\).  Thus
\(\mathcal G_m\subseteq\mathcal F_m\).  The hidden clauses are already
satisfied by their root literal at \(t\), so their nonroot parts do not
affect \(S_m\).  In particular \(S_m,J_m,s_{j,m},v_m\), and
\(\mathcal W_m\) are \(\mathcal G_m\)-measurable.  Conditional on
\(\mathcal G_m\), the hidden nonroot parts are independent uniform
proper two-clauses on the \(u\) nonroots.  Indeed, conditioning the
independent clauses on their types separately preserves their product
law, and revealing nonroot parts at other indices does not change it.

Put
\[
 d=\lceil\log n\rceil,\qquad q_m=q_d(S_m).
\]

\begin{lemma}[Root-flip tests under prefix conditioning]
\label{lem:prefix-root-tests}
For every \(i\le m<L\), on the \(\mathcal G_m\)-measurable event
\(\{\max_{j\in R}s_{j,m}\le d,\ |J_m|\ge1\}\),
\begin{equation}\label{eq:prefix-root-tests}
 \PP(A_m\mid\mathcal G_m)
 \le\ind{S_m\ne\varnothing}q_m^{|J_m|}.
\end{equation}
\end{lemma}

\begin{proof}
If \(A_m\) occurs, then for each \(j\in J_m\) the pair clauses at
indices in \(\mathcal I_{j,m}\) must kill \(S_m\).  Otherwise choose
an assignment in \(S_m\) satisfying those pairs, extend it with root
values \(t\), and flip only \(j\).  A clause with no true root literal
before the flip has a satisfied nonroot part, which is unchanged.  A
clause with a true root literal at a root other than \(j\) remains
satisfied.  Any remaining clause has \(j\) as its unique true root
before the flip.  It cannot be a collision because \(j\in J_m\), so
it belongs to \(\mathcal I_{j,m}\) and its pair is satisfied by the
chosen assignment.  The flipped assignment therefore satisfies
\(B_m\) with root values different from \(t\), contradicting
\eqref{eq:prefix-delay-events}.

The required killing events for different \(j\) use disjoint families
of the conditionally independent hidden pairs.  On a nonempty \(S_m\),
each family has at most \(d\) members and therefore kills with
probability at most \(q_d(S_m)=q_m\).  A family of zero pairs has
killing probability zero, so it causes no exception.  Multiplication of
these conditional bounds, together with the necessity of
\(S_m\ne\varnothing\), proves \eqref{eq:prefix-root-tests}.
\end{proof}

\subsection{Truncating incidence counts and summing occupation}

It remains to sum the right side of \eqref{eq:prefix-weighted-delay}.
The tests make deletion unlikely when \(q_m\) is small.  When it is
large, the lifetime estimate limits both the remaining deletion time
and the number of indices at that level.  We first isolate the rare
prefixes for which too many tests or collisions are present.

\begin{lemma}[Prefix incidence counts]\label{lem:prefix-counts}
Uniformly in \(i\le m<L\) and in the fixed value of \(C_i\),
\begin{align}
 \PP\bigl(\max_{j\in R}s_{j,m}>d\bigr)&=O(n^{-4}),
 \label{eq:prefix-large-count}\\
 \PP(v_m\ge1)&=O(n^{-1}),\qquad
 \PP(v_m\ge2)=O(n^{-2}).
 \label{eq:prefix-collision-counts}
\end{align}
\end{lemma}

\begin{proof}
A baseline clause contains a specified root with probability \(3/n\),
independently across indices, and \(s_{j,m}\) is no greater than that
root's incidence count.  A union bound over roots and choices of \(d\)
indices among at most \(L=10n\) trials gives
\begin{equation}\label{eq:prefix-count-numerical}
 \PP\bigl(\max_{j\in R}s_{j,m}>d\bigr)
 \le3\binom Ld(3/n)^d\le3(30e/d)^d=O(n^{-4}).
\end{equation}
The final bound holds for sufficiently large \(n\), since
\(d=\lceil\log n\rceil\) and \(\log d\to\infty\).

A trial contains any specified pair of roots with probability
\(6/[n(n-1)]\).  The union over the three pairs bounds its collision
probability by \(18/[n(n-1)]\).  A union bound over at most \(L\)
independent trials gives \(\PP(v_m\ge1)=O(n^{-1})\); a union bound
over pairs of trials gives
\(\PP(v_m\ge2)\le\binom L2 O(n^{-4})=O(n^{-2})\).
\end{proof}

With \(K=64\), set
\begin{equation}\label{eq:prefix-cutoff}
 a_0=\frac{4Kd}{u^2}=\frac{256d}{(n-3)^2}.
\end{equation}
For sufficiently large \(n\), \(a_0<1\).  On \(S_m\ne\varnothing\),
\begin{equation}\label{eq:prefix-weighted-lifetime}
 \mathcal W_m q_m^{3/2}\le C d^{3/2}.
\end{equation}
For \(q_m\ge a_0\), this follows by applying
Proposition~\ref{prop:prefix-lifetime} with \(a=q_m\), and then using
\(\mathcal W_m\le\tau(S_m)/p_0\).  For \(q_m<a_0\), including
\(q_m=0\), use \(\mathcal W_m\le L\) and
\begin{equation}\label{eq:prefix-small-level}
 L a_0^{3/2}
 =10n\left(\frac{256d}{(n-3)^2}\right)^{3/2}
 =O(d^{3/2}n^{-2})\le C d^{3/2}.
\end{equation}

\begin{lemma}[Reduction to occupation]\label{lem:prefix-reduction}
For all sufficiently large \(n\), uniformly in \(i\) and in the value
of \(C_i\),
\begin{equation}\label{eq:prefix-reduction}
 \EE D_i^2\le C(1+d^{3/2})+
 C d^{3/2}\sum_{m=i}^{L-1}
       \EE\bigl[\ind{S_m\ne\varnothing}q_m^{3/2}\bigr].
\end{equation}
\end{lemma}

\begin{proof}
In \eqref{eq:prefix-weighted-delay}, the terms on
\(\{\max_j s_{j,m}>d\}\cup\{v_m\ge2\}\) contribute at most
\[
 L\sum_{m=i}^{L-1}
 \bigl[\PP(\max_j s_{j,m}>d)+\PP(v_m\ge2)\bigr]
 =O\bigl(L^2(n^{-4}+n^{-2})\bigr)=O(1).
\]
For every remaining term, the restrictions, \(\mathcal W_m\), and
\(S_m\) are \(\mathcal G_m\)-measurable.  If \(v_m=1\), then
\(|J_m|\ge2\).  Conditioning on \(\mathcal G_m\) in
Lemma~\ref{lem:prefix-root-tests} bounds its weighted probability by
\(\ind{S_m\ne\varnothing}\mathcal W_mq_m^2\).  On a nonempty set,
\eqref{eq:prefix-weighted-lifetime} and \(q_m\le1\) bound this by
\(Cd^{3/2}\).  The one-collision contribution is therefore at most
\[
 Cd^{3/2}\sum_{m=i}^{L-1}\PP(v_m=1)=O(d^{3/2}).
\]
No independence between \(v_m\) and the surviving set is asserted or
needed.

If \(v_m=0\), all three roots are available, and the same conditional
bound has exponent three.  By \eqref{eq:prefix-weighted-lifetime},
\[
 \ind{S_m\ne\varnothing}\mathcal W_mq_m^3
 \le Cd^{3/2}\ind{S_m\ne\varnothing}q_m^{3/2}.
\]
We may now discard the incidence restrictions.  Summing the three
cases and absorbing the factor three from
\eqref{eq:prefix-weighted-delay} proves \eqref{eq:prefix-reduction}.
\end{proof}

We next bound how long the decreasing sets \(S_m\) can remain nonempty
at a given pair-killing level.  The first visit to that level determines
a set whose future lifetime controls all subsequent visits.

\begin{lemma}[Occupation above a level]\label{lem:prefix-occupation}
For all sufficiently large \(n\) and every \(a_0\le a\le1\),
\begin{equation}\label{eq:prefix-occupation}
 \sum_{m=i}^{L-1}\PP(S_m\ne\varnothing,\ q_m\ge a)
 \le Cd^{3/2}a^{-3/2}.
\end{equation}
For \(0<a<a_0\), the same sum is at most \(L\).
\end{lemma}

\begin{proof}
Let \(\sigma_a\) be the first index in \(\{i,\ldots,L-1\}\) for
which \(S_m\ne\varnothing\) and \(q_m\ge a\), and set
\(\sigma_a=\infty\) when none exists.  This is a stopping time for
\(\mathcal F_m\), because \(q_d\) is a deterministic function of
its set argument.  On \(\{\sigma_a=m\}\), every counted index is
at least \(m\).  Also \(S_\ell=\varnothing\) for
\(\ell\ge m+U_m\): by then the future ordinary clauses have killed
\(S_m\), and any member of \(S_\ell\) would be a member of \(S_m\)
satisfying all those clauses.  The number of counted indices is thus
at most \(U_m\), including the first index \(m\).

For \(a\ge a_0\), the event \(\{\sigma_a=m\}\) is
\(\mathcal F_m\)-measurable, and on it the lifetime proposition applies.
The fresh-future identity gives
\begin{align*}
 \EE[\ind{\sigma_a=m}U_m]
 &=\EE\left[\ind{\sigma_a=m}\frac{\tau(S_m)}{p_0}\right]\\
 &\le Cd^{3/2}a^{-3/2}\PP(\sigma_a=m).
\end{align*}
Summing over the disjoint first-visit events proves
\eqref{eq:prefix-occupation}.  This conditioning is legitimate because
\(m\ge i\) and all trials after \(m\) remain unexposed.  At a smaller
level there are at most \(L\) indices to count.
\end{proof}

\begin{proof}[Proof of Theorem~\ref{thm:prefix-variance}]
For \(0\le q\le1\),
\(q^{3/2}=\int_0^1\frac32a^{1/2}\ind{q\ge a}\,da\).
Applying this identity and Lemma~\ref{lem:prefix-occupation} gives
\begin{align*}
 \sum_{m=i}^{L-1}\EE[\ind{S_m\ne\varnothing}q_m^{3/2}]
 &=\int_0^1\frac32a^{1/2}
       \sum_{m=i}^{L-1}\PP(S_m\ne\varnothing,q_m\ge a)\,da\\
 &\le L a_0^{3/2}+Cd^{3/2}\int_{a_0}^1\frac{da}{a}
 \le Cd^{3/2}\log n.
\end{align*}
The final step uses \eqref{eq:prefix-small-level} and
\(\log(1/a_0)\le2\log n\), since \(a_0\ge n^{-2}\).  The sum is finite
and the integrands are
nonnegative, so the interchange of sum, expectation, and integral is
valid.  Lemma~\ref{lem:prefix-reduction} now yields, uniformly in the
deleted index and the fixed clause,
\begin{equation}\label{eq:prefix-squared-influence}
 \EE D_i^2\le C(1+d^{3/2})+Cd^3\log n\le C(\log n)^4.
\end{equation}
Average over \(C_i\) to remove the conditioning, and sum
\eqref{eq:prefix-squared-influence} in \eqref{eq:prefix-omission}.
This gives \(\Var(T_n)\le Cn(\log n)^4\) for sufficiently large
\(n\).  Enlarging \(C\) covers the remaining integers \(n\ge3\),
since \(0\le T_n\le10n\).  Chebyshev's inequality proves the stated
tail bound.
\end{proof}

For every fixed \(0<\eta<1/2\), the \(\eta\) and \(1-\eta\)
quantiles of \(T_n\) therefore lie within
\(O_\eta(\sqrt n\log^2 n)\) of \(m_n\).  The probability outside
any diverging multiple of this scale tends to zero.  These statements
use the finite-size mean \(m_n\).

\subsection{Cap changes and the center comparison}

The prefix proof supplies a second concentration input for
Theorem~\ref{thm:proper-center-module}.  That theorem assumes a
concentration bound and proves the center comparisons independently of
how the bound is obtained.  We apply it here using the variance bound of
Theorem~\ref{thm:prefix-variance}, together with the cap identities
already proved in Section~\ref{sec:three-clause-refinements}.  Put
\[
 W_n=\min\{H_n-1,10n\},\qquad V_n=\min\{H_n-1,16n\},
 \qquad e_n=\frac{\EE T_n}{n},\qquad P_n(j)=\PP(H_n>j).
\]

\begin{corollary}\label{cor:prefix-consequences}
The variance and tail bounds of Theorem~\ref{thm:prefix-variance} hold
also for \(W_n\) and \(V_n\), each centered at its own expectation.
Moreover \(e_n\to\alpha_3\), the strict-side threshold of the proper
three-clause model: for every fixed \(c\ge0\) with \(c\ne\alpha_3\),
\[
 P_n(\lfloor cn\rfloor)\longrightarrow
 \begin{cases}1,&c<\alpha_3,\\0,&c>\alpha_3.\end{cases}
\]
There is an absolute constant \(C\) such that for all sufficiently large
\(n=\min\{n_1,n_2\}\) with \(\max\{n_1,n_2\}\le3n\), and for all
sufficiently large \(n\) in the second inequality,
\begin{equation}\label{eq:prefix-center-comparisons}
 e_{n_1+n_2}\ge\min\{e_{n_1},e_{n_2}\}-Cn^{-1/4},
 \qquad e_{n+1}\ge e_n-Cn^{-1/4}.
\end{equation}
\end{corollary}

\begin{proof}
Set \(\rho=2(7/8)^{10}<1\).  With \(T=T_n\), \(W=W_n\), and
\(V=V_n\), identities~\eqref{eq:three-same-cap}--\eqref{eq:three-cross-cap}
show that the mean corrections from \(m_n-1\) are at most \(\rho^n\)
and \((6n+1)\rho^n\), respectively.  The corresponding absolute differences
of standard deviations are at most \(\rho^{n/2}\) and
\((6n+1)\rho^{n/2}\); the sharper geometric bounds are those in
Section~\ref{sec:three-clause-refinements}, obtained from
Proposition~\ref{prop:cap-bridge}.  These exponentially small corrections
transfer the variance estimate to both last-satisfiable statistics.
Chebyshev's inequality gives their tails about their own means.

For the center claim, choose a fixed \(0<c<e^{-2}\) and put
\(j=\lfloor cn\rfloor\).  Lemma~\ref{lem:sparse-satisfiability} gives
\(\EE T_n\ge jP_n(j)=cn-o(n)\).  The survival-sum bound
\eqref{cmp:mean-upper} with \(k=3\) and density six gives
\(\EE T_n\le\EE H_n\le6n+8[2(7/8)^6]^n\).
Thus \(c_0\le e_n\le7<10\) eventually, for some absolute \(c_0>0\).
The relevant survival identity is
\(P_n(j)=\PP(T_n>j)\) for integers \(0\le j<10n\);
the restriction below the cap is essential, whereas the last-SAT identity
\(P_n(j)=\PP(W_n\ge j)\) includes \(j=10n\).
Theorem~\ref{thm:prefix-variance} supplies the concentration input
\begin{equation}\label{eq:prefix-center-concentration}
 \PP\bigl(|T_n-m_n|\ge n^{3/4}\bigr)\le\eta_n,
 \qquad \eta_n=C\frac{(\log n)^4}{\sqrt n}\longrightarrow0.
\end{equation}
Apply Theorem~\ref{thm:proper-center-module} with \(k=3\), \(B=10\),
\(a=c_0\), \(A=7\), and \(\delta=1/4\).  It gives both comparisons
in~\eqref{eq:prefix-center-comparisons}, in particular for the stated
comparable sizes, and convergence with the proper strict-side limits.
Uniqueness identifies the limit with \(\alpha_3\) of
\cite[Theorem~\ThresholdTheorem]{FixedKThreshold}.  The one-sided
estimate~\eqref{cmp:tree-bound} also gives, for every sufficiently large
\(s\) and every \(N\ge s^2\),
\[
 e_N\ge e_s-C's^{-1/4}.
\]
\end{proof}

For the comparable-size consequence, the coarser trial comparison gives
a concrete error bound.  With \(N=n_1+n_2\),
\(n=\min(n_1,n_2)\), \(e_*=\min(e_{n_1},e_{n_2})\), set
\(\varepsilon=n^{-1/4}\) and
\(M=\lfloor(e_*-2\varepsilon)N\rfloor\), as in the proof of
Theorem~\ref{thm:proper-center-module}.  The anchors give
\((c_0/4)N\le M\le7N\) for large \(n\).  Since \(C_3=27\) in
Lemma~\ref{cmp:trial-comparison} and \(N\le4n\),
\eqref{cmp:coarse-replacement} bounds the loss by
\begin{equation}\label{eq:prefix-756-loss}
 \frac{27M}{\varepsilon^3n^3}
 \le\frac{27\cdot7\cdot4n}{n^{9/4}}
 =756n^{-5/4}.
\end{equation}
The restriction \(\max(n_1,n_2)\le3n\) is used only for this numerical
bound.  For neighboring sizes, set \(N=n+1\) and
\(M=\lfloor m_n-2n^{3/4}\rfloor\).  The constant \(D_3=729/8\) gives loss
\((729/8)M/(\varepsilon^3n^3)=O(n^{-5/4})\).
In either case, the count-transfer index
\(j=\lfloor(1-\varepsilon)M-N^{3/4}\rfloor\) is nonnegative for
large \(n\) and lies strictly below \(10N\), as required in
\eqref{cmp:count-transfer}.

The deviation bounds above are about each statistic's finite-size mean.
The comparison inequalities are one-sided and do not give a two-sided
rate for \(e_n-\alpha_3\).
